\documentclass[12pt,a4paper]{amsart}

\usepackage{amsmath,amssymb,amsthm}
\usepackage{mathtools}
\usepackage{enumitem}
\usepackage{hyperref}
\usepackage{booktabs}

%%%─── Theorem environments
\newtheorem{theorem}{Theorem}[section]
\newtheorem{lemma}[theorem]{Lemma}
\newtheorem{proposition}[theorem]{Proposition}
\newtheorem{corollary}[theorem]{Corollary}
\newtheorem{conjecture}[theorem]{Conjecture}
\newtheorem{hypothesis}[theorem]{Hypothesis}
\theoremstyle{definition}
\newtheorem{definition}[theorem]{Definition}
\newtheorem{assumption}[theorem]{Assumption}
\newtheorem{problem}[theorem]{Open Problem}
\theoremstyle{remark}
\newtheorem{remark}[theorem]{Remark}
\newtheorem*{remark*}{Remark}

%%%─── Macros (consistent with series)
\newcommand{\norm}[1]{\|#1\|}
\newcommand{\ip}[2]{\langle #1,\, #2\rangle}
\newcommand{\Hlim}{\mathcal{H}_{\lim}}
\newcommand{\Hc}{\mathcal{H}_c}
\newcommand{\Phiinf}[1]{\Phi_{#1}^{(\infty)}}
\newcommand{\Phic}[2]{\Phi_{#2}^{(#1)}}
\newcommand{\lambdainf}[1]{\lambda_{#1}^{(\infty)}}
\newcommand{\lambdac}[2]{\lambda_{#2}^{(#1)}}
\newcommand{\PW}{PW}
\newcommand{\spc}[1]{\sigma(#1)}
\newcommand{\Kspace}{\mathcal{K}}

\subjclass[2020]{47A10, 47B25, 42C05, 46B15, 46C05}
\keywords{Prolate spheroidal wave functions, completeness,
  orthonormal basis, strong operator convergence,
  Paley--Wiener space, structural classification,
  Hilbert--P\'{o}lya programme}

\begin{document}

\title[Completeness of the Limiting Slepian Eigenbasis]{%
  Completeness of the Limiting Slepian Eigenbasis in $L^2(\mathbb{R})$:
  A Structural Classification}

\author{[Author]}
\address{[Address]}
\email{[Email]}

\begin{abstract}
Let $\Hlim$ denote the bounded self-adjoint SOT-limit operator
constructed in Paper~X, and let $\{\Phiinf{n}\}_{n \ge 0}$
be the orthonormal system of limit eigenfunctions from Paper~XI.
Define $\Kspace := \overline{\mathrm{span}}\{\Phiinf{n}\} \subseteq L^2(\mathbb{R})$.
This paper studies the completeness problem $\Kspace = L^2(\mathbb{R})$,
which is Open Problem~7 of Paper~IX and Open Problem~5 of Paper~X.

\medskip
\textbf{Main contribution: structural classification.}
We classify the completeness problem into three mechanisms
of fundamentally different logical character:

\medskip
\begin{center}
\renewcommand{\arraystretch}{1.3}
\begin{tabular}{lll}
\toprule
Mechanism & Character & Status \\
\midrule
A (Frame identity) & \emph{Equivalence reformulation}
  & Not an independent proof route \\
B (Paley--Wiener transfer) & \emph{Genuine analytic reduction}
  & Conditional on Hyp.~\ref{hyp:coeff-stability}+\ref{hyp:reconstruction} \\
C (Compact resolvent) & \emph{Structural collapse}
  & Conditional on Hyp.~\ref{hyp:compact-resolvent} \\
\bottomrule
\end{tabular}
\end{center}

\medskip
The central result (Theorem~\ref{thm:route-A-equivalence}) establishes
that the Parseval identity $\sum_n |\ip{f}{\Phiinf{n}}|^2 = \norm{f}^2$
is \emph{equivalent} to $\Kspace = L^2(\mathbb{R})$; it is a structural
characterisation, not an independent route.
This shows that frame identities in this setting are not global
proof tools but characterisations of the global closure of the
spectral representation.

\medskip
\textbf{Main theorem (Route~B).}
Route~B reduces completeness to two independent hypotheses:
(H1) coefficient stability under SOT-limit
(grounded in Papers~VIII and XII)
and (H2) reconstruction of $\PW_c^2$-elements
in the limit basis (the genuine open step).
Route~B is not circular: it establishes that
$\overline{\Kspace}$ contains each $\PW_c^2$
--- a local completeness property in the inductive
limit of Paley--Wiener spaces --- and closes via density.

\medskip
\textbf{Conditional result (Route~C).}
Under Hypothesis~\ref{hyp:compact-resolvent},
$\{\Phiinf{n}\}$ is an orthonormal basis;
the spectrum of $\Hlim$ is purely discrete with
finite multiplicities, accumulating only at $0$ or $1$.
This hypothesis is mutually exclusive with
$\spc{\Hlim} = [0,1]$ (Paper~X, Open Problem~4).
\end{abstract}

\maketitle

%%═══ 1. INTRODUCTION
\section{Introduction}
\label{sec:intro}

\subsection{Setting and goal}

\begin{center}
\fbox{\parbox{0.88\linewidth}{%
\centering\itshape
Frame identities in this setting are not global proof tools,
but characterisations of the global closure of the spectral
representation. Any frame-based argument must therefore be
localised to the level of individual Paley--Wiener spaces.}}
\end{center}

\medskip
\noindent
The Slepian concentration operator $\Hc$ acts on $L^2(\mathbb{R})$ via
\[
  (\Hc f)(x) = \int_{-\infty}^{\infty}
    \frac{\sin c(x-y)}{\pi(x-y)} f(y)\, dy.
\]
Paper~X establishes $\Hlim := \mathrm{SOT}\text{-}\lim_{c \to \infty} \Hc$
as a bounded self-adjoint operator with $\norm{\Hlim} \le 1$
and $\spc{\Hlim} \subseteq [0,1]$.
Paper~XI constructs eigenfunctions $\{\Phiinf{n}\}$ satisfying
\begin{equation}\label{eq:eigen-relation}
  \Hlim \Phiinf{n} = \lambdainf{n} \Phiinf{n}, \qquad
  \ip{\Phiinf{m}}{\Phiinf{n}} = \delta_{mn},
  \qquad \lambdainf{n} \in (0,1).
\end{equation}
The central question is whether
\begin{equation}\label{eq:completeness-question}
  \Kspace := \overline{\mathrm{span}}\{\Phiinf{n}\}_{n \ge 0}
  \stackrel{?}{=} L^2(\mathbb{R}).
\end{equation}

\subsection{What this paper does not claim}

\begin{itemize}
\item We do not claim $\spc{\Hlim}$ is purely discrete
  (Paper~X, Open Problem~4 remains open).
\item We do not derive completeness from Sturm--Liouville
  theory of $\mathcal{L}_c$: the operators $\mathcal{L}_c$
  and $\Hc$ are distinct; their eigenfunctions coincide
  classically \cite{Osipov2013} but the spectral
  structure of $\Hlim$ is independent.
\item We do not re-prove self-adjointness of $\Hlim$
  (unconditional in Paper~X).
\end{itemize}

\subsection{Role in the series}

\begin{center}
\begin{tabular}{lll}
\toprule
Paper & Content & Status \\
\midrule
X  & SOT: $\Hc \xrightarrow{\mathrm{SOT}} \Hlim$;
     Mosco; Friedrichs & Final \\
XI & System $\{\Phiinf{n}\}$; orthonormality;
     spectral structure & Final \\
XII & Localization principle; PSWF obstruction set
     & Final (submitted) \\
\textbf{XIII} & \textbf{Structural classification
     of completeness} & \textbf{This paper} \\
\bottomrule
\end{tabular}
\end{center}

%%═══ 2. INHERITED DATA
\section{Inherited Data}
\label{sec:input}

\begin{lemma}[Series results used in this paper]
\label{lem:inherited}
\begin{enumerate}[\rm(A)]
\item\label{A:bounded-sa}
  \emph{(Paper~X, Thm.~A.)} $\Hlim$ is bounded self-adjoint,
  $\norm{\Hlim} \le 1$, $\spc{\Hlim} \subseteq [0,1]$.
  In particular $(\Hlim + I)^{-1}$ is a well-defined bounded operator.

\item\label{A:ONS}
  \emph{(Paper~XI, Lem.~3.1 + Thm.~5.1.)}
  $\{\Phiinf{n}\}_{n \ge 0}$ satisfies \eqref{eq:eigen-relation}
  and is orthonormal.

\item\label{A:nonvanish}
  \emph{(Paper~VIII, Thm.~\textup{thm:nonvanish}.)}
  For each $n \le \kappa c$, there exists
  $E_c \subset \{|t - t_*(c,n)| \le c^{-1/2}\}$
  with $|E_c| \ge c_* c^{-1/2}$ and
  $|\widehat{\Phic{c}{n}}(t)| \ge c_1 c^{-1/2}$ on $E_c$.
  The constants $c_*, c_1$ depend only on $\kappa$.

\item\label{A:precompact}
  \emph{(Paper~IX, Lem.~3.1.)}
  $\{\Phic{c}{n}\}_{c > 0}$ is precompact in $L^2(\mathbb{R})$,
  via Fr\'echet--Kolmogorov: uniform exponential decay
  and uniform $H^1$-bound \cite{Osipov2013}.

\item\label{A:eigenvalue-rate}
  \emph{(Paper~VIII, Cor.~4.3.)}
  $|\lambdac{c}{n} - \lambdainf{n}| \le C_\kappa c^{-1/4}$
  for $n \le \kappa c$, uniformly in $n$.

\item\label{A:eigenvec-conv}
  \emph{(Paper~XII, conditional on Gap-S + SOT-Faster.)}
  Under those two hypotheses,
  $\norm{\Phic{c}{n} - \Phiinf{n}} \to 0$
  as $c \to \infty$ for each fixed $n$.
\end{enumerate}
\end{lemma}

\begin{remark}[SOT does not control spectral type]
\label{rem:sot-limitation}
SOT-convergence $\Hc \xrightarrow{\mathrm{SOT}} \Hlim$ does not imply
norm-resolvent convergence and does not determine whether
$\spc{\Hlim}$ is discrete, continuous, or equal to $[0,1]$
\cite[Ex.~VIII.7.14]{ReedSimon1975}.
\end{remark}

%%═══ 3. INTERFACE LEMMA
\section{Interface Lemma: Completeness via Orthogonality}
\label{sec:interface}

\begin{lemma}[Orthogonality characterisation]
\label{lem:orth-char}
For an orthonormal system $\{\Phiinf{n}\}$ in $L^2(\mathbb{R})$:
\[
  \Kspace = L^2(\mathbb{R})
  \quad\Longleftrightarrow\quad
  \Bigl[\,\ip{f}{\Phiinf{n}} = 0\ \forall\, n
  \;\Rightarrow\; f = 0\,\Bigr].
\]
\end{lemma}
\begin{proof}
\cite[Theorem~3.3]{Conway1990}.
\end{proof}

\begin{remark}[Function of this lemma]
This is the \emph{interface lemma}: it converts the
geometric question~\eqref{eq:completeness-question}
into an analytic condition on inner products,
which is the entry point for Routes~B and~C.
\end{remark}

%%═══ 4. ROUTE A: FRAME IDENTITY AS EQUIVALENCE
\section{Route~A: The Frame Identity is an Equivalence, Not a Reduction}
\label{sec:riesz}

We record Route~A in order to establish its precise logical status.

\begin{definition}[Gram operator]
\label{def:gram}
The Gram operator $G \colon L^2(\mathbb{R}) \to L^2(\mathbb{R})$ is
\[
  Gf := \sum_{n \ge 0} \ip{f}{\Phiinf{n}} \Phiinf{n}.
\]
Bessel's inequality gives $\norm{Gf} \le \norm{f}$,
so $G$ is a well-defined bounded contraction.
$G$ is the orthogonal projection onto $\Kspace$.
\end{definition}

\begin{theorem}[Frame identity $\Longleftrightarrow$ completeness]
\label{thm:route-A-equivalence}
For the orthonormal system $\{\Phiinf{n}\}$,
the following are equivalent:
\begin{enumerate}[\rm(i)]
\item $\Kspace = L^2(\mathbb{R})$ (completeness);
\item $G = \mathrm{id}_{L^2(\mathbb{R})}$;
\item $\displaystyle\sum_{n \ge 0}|\ip{f}{\Phiinf{n}}|^2 = \norm{f}^2$
  for all $f \in L^2(\mathbb{R})$ (Parseval identity);
\item There exists $A > 0$ such that
  $\displaystyle\sum_{n \ge 0}|\ip{f}{\Phiinf{n}}|^2 \ge A\norm{f}^2$
  for all $f \in L^2(\mathbb{R})$ (frame lower bound).
\end{enumerate}
\end{theorem}
\begin{proof}
\textbf{(i)$\Rightarrow$(ii).}
If $\Kspace = L^2(\mathbb{R})$, then $G$ projects onto all
of $L^2(\mathbb{R})$, so $G = \mathrm{id}$.

\textbf{(ii)$\Rightarrow$(iii).}
For any $f \in L^2(\mathbb{R})$:
$\norm{Gf}^2 = \sum_n |\ip{f}{\Phiinf{n}}|^2$.
If $G = \mathrm{id}$, this equals $\norm{f}^2$.

\textbf{(iii)$\Rightarrow$(iv).}
Take $A = 1$.

\textbf{(iv)$\Rightarrow$(i).}
Suppose $\sum_n |\ip{f}{\Phiinf{n}}|^2 \ge A\norm{f}^2 > 0$
for all $f \ne 0$.
If $f \perp \Phiinf{n}$ for all $n$,
then the left-hand side is $0$, forcing $f = 0$.
Hence $\{\Phiinf{n}\}^\perp = \{0\}$,
which is equivalent to $\Kspace = L^2(\mathbb{R})$
by Lemma~\ref{lem:orth-char}.

\textbf{Remark on the value of $A$.}
For orthonormal systems, Bessel's inequality gives
$\sum_n |\ip{f}{\Phiinf{n}}|^2 \le \norm{f}^2$,
so if (iv) holds then necessarily $A = 1$
and (iii) holds, confirming (iii)$\Leftrightarrow$(iv).
\end{proof}

\begin{remark}[Route~A: structural characterisation only]
\label{rem:route-A-status}
Theorem~\ref{thm:route-A-equivalence} shows that
any proof of the frame lower bound
$\sum_n |\ip{f}{\Phiinf{n}}|^2 \ge A\norm{f}^2$
is precisely as hard as proving completeness itself.
Hence:
\begin{itemize}
\item Frame identities in this setting are not global
  proof tools, but characterisations of the global
  closure of the spectral representation.
\item The nonvanishing estimates of Paper~VIII
  (Lemma~\ref{lem:inherited}\ref{A:nonvanish})
  control $|\widehat{\Phi}_n^{(c)}|$ for
  \emph{finite} $c$ only.
  Passing to the limit in the frame sum requires
  coefficient stability, which is exactly the content
  of Hypotheses~\ref{hyp:coeff-stability}--\ref{hyp:reconstruction}.
\item Any frame-based argument must be
  \emph{localised to the level of individual
  Paley--Wiener spaces} to avoid circularity.
\end{itemize}
\end{remark}

\begin{problem}[Structural decomposition of completeness]
\label{prob:frame-stability}
Analyse whether $\Kspace = L^2(\mathbb{R})$ decomposes
into the following two logically independent statements:
\begin{enumerate}[\rm(a)]
\item \emph{(Limit-interchange stability)}
\[
  \lim_{c \to \infty}
  \sum_{n \le \kappa c}
  |\ip{f}{\Phic{c}{n}}|^2
  = \sum_{n \ge 0}|\ip{f}{\Phiinf{n}}|^2
  \quad \forall\, f \in L^2(\mathbb{R}).
\]
\item \emph{(Frame isometry)}
$\displaystyle\sum_{n \ge 0}|\ip{f}{\Phiinf{n}}|^2 = \norm{f}^2$.
\end{enumerate}
By Theorem~\ref{thm:route-A-equivalence}, statement~(b) is
precisely the completeness problem.
Statement~(a) is a necessary consequence of completeness
(not a sufficient one); it follows from dominated convergence
and strong eigenvector convergence
(Lemma~\ref{lem:inherited}\ref{A:eigenvec-conv}).
\end{problem}

%%═══ 5. ROUTE B: PALEY–WIENER TRANSFER (MAIN THEOREM)
\section{Route~B: Paley--Wiener Transfer --- Main Theorem}
\label{sec:pw}

Route~B avoids circularity by working locally:
it establishes that $\Kspace$ contains each $\PW_c^2$
--- a local completeness property in the inductive
limit of Paley--Wiener spaces ---
then closes via density of $\bigcup_{c>0}\PW_c^2$.

\begin{lemma}[Paley--Wiener density]
\label{lem:pw-dense}
$\bigcup_{c > 0} \PW_c^2$ is dense in $L^2(\mathbb{R})$.
\end{lemma}
\begin{proof}
$\bigcup_{c > 0} \PW_c^2$ contains all Schwartz functions
with compactly supported Fourier transform,
a dense subspace \cite[Theorem~7.14]{Rudin1991}.
\end{proof}

\medskip
Hypothesis~\ref{hyp:pw-transfer} from the previous version
is now split into two independent hypotheses,
corresponding to the two independent open steps.

\begin{hypothesis}[Coefficient stability, H1]
\label{hyp:coeff-stability}
For each fixed $c_0 > 0$, $g \in \PW_{c_0}^2$,
and $c \ge c_0$, let
$a_n^{(c)} := \ip{g}{\Phic{c}{n}}$.
As $c \to \infty$, the coefficients converge
in $\ell^2$:
\[
  a_n^{(c)} \longrightarrow a_n^{(\infty)}
  \quad \text{in } \ell^2(\mathbb{N}),
  \qquad
  a_n^{(\infty)} := \ip{g}{\Phiinf{n}}.
\]
\end{hypothesis}

\begin{remark}[Grounding of Hypothesis~\ref{hyp:coeff-stability}]
\label{rem:H1-grounding}
For each fixed $n$, the pointwise convergence
$a_n^{(c)} \to a_n^{(\infty)}$ follows from strong
eigenvector convergence
(Lemma~\ref{lem:inherited}\ref{A:eigenvec-conv})
conditionally on Gap-S and SOT-Faster (Paper~XII).
Upgrading to $\ell^2$-convergence requires uniform
tail control: $\sum_{n > N} |a_n^{(c)}|^2 \to 0$
uniformly in $c$, which follows from the Bessel majorant
$\sum_n |a_n^{(c)}|^2 \le \norm{g}^2$
and dominated convergence.
The non-degeneracy of individual $a_n^{(c)}$ for
$n \le \kappa c$ is supported by
Lemma~\ref{lem:inherited}\ref{A:nonvanish}.
\end{remark}

\begin{hypothesis}[Reconstruction in $\PW_c^2$, H2]
\label{hyp:reconstruction}
Under the notation of Hypothesis~\ref{hyp:coeff-stability},
for each $c_0 > 0$ and $g \in \PW_{c_0}^2$:
\[
  g = \sum_{n \ge 0} a_n^{(\infty)} \Phiinf{n}
  \quad \text{in } L^2(\mathbb{R}).
\]
In other words, $\overline{\Kspace}$ contains
every element of $\PW_{c_0}^2$.
\end{hypothesis}

\begin{remark}[H2 is the genuine open step]
\label{rem:H2-status}
Hypothesis~\ref{hyp:reconstruction} is the substantive
content of Route~B.
It asserts a \emph{local completeness property}:
$\PW_{c_0}^2 \subseteq \Kspace$ for each fixed $c_0$.
This is structurally weaker than global completeness
$\Kspace = L^2(\mathbb{R})$, but it is still a local
completeness property in the inductive limit of
Paley--Wiener spaces.
H2 does not follow from H1 alone;
it requires that the limit-coefficient expansion
converges \emph{to the original function $g$},
not merely in an abstract $\ell^2$-sense.
H1 provides the coordinate control; H2 requires
that these coordinates reconstruct $g$ in $L^2$.
\end{remark}

\begin{theorem}[Main theorem: Route~B]
\label{thm:main-B}
\textup{(Conditional on Hypotheses~\ref{hyp:coeff-stability}
and~\ref{hyp:reconstruction}.)}
If both hypotheses hold,
then $\Kspace = L^2(\mathbb{R})$.
\end{theorem}
\begin{proof}
Suppose $f \in L^2(\mathbb{R})$ with
$\ip{f}{\Phiinf{n}} = 0$ for all $n \ge 0$.

\textbf{Step~1.}
Fix $c_0 > 0$ and $g \in \PW_{c_0}^2$.
By Hypothesis~\ref{hyp:reconstruction},
$g = \sum_{n \ge 0} a_n^{(\infty)} \Phiinf{n}$
in $L^2(\mathbb{R})$.

\textbf{Step~2.}
Compute:
\[
  \ip{f}{g}
  = \ip*{f}{\,\sum_{n \ge 0} a_n^{(\infty)} \Phiinf{n}}
  = \sum_{n \ge 0} a_n^{(\infty)} \ip{f}{\Phiinf{n}}.
\]
The interchange is justified: by Hypothesis~\ref{hyp:coeff-stability},
$(a_n^{(\infty)}) \in \ell^2$, and by Bessel's inequality
$(\ip{f}{\Phiinf{n}}) \in \ell^2$, so the series converges
absolutely by Cauchy--Schwarz.

\textbf{Step~3.}
Since $\ip{f}{\Phiinf{n}} = 0$ for all $n$, we get
$\ip{f}{g} = 0$.

\textbf{Step~4.}
As $c_0 > 0$ and $g \in \PW_{c_0}^2$ were arbitrary,
$f \perp \bigcup_{c_0 > 0} \PW_{c_0}^2$.
By density (Lemma~\ref{lem:pw-dense}), $f = 0$.

\textbf{Conclusion.}
By Lemma~\ref{lem:orth-char}, $\Kspace = L^2(\mathbb{R})$.
\end{proof}

\begin{remark}[Why Route~B is not circular]
\label{rem:route-B-not-circular}
The argument requires that $\Kspace$ contains each
$\PW_{c_0}^2$ (Hypothesis~\ref{hyp:reconstruction}),
not that $\Kspace = L^2(\mathbb{R})$ a priori.
More precisely: the argument establishes
\[
  \PW_{c_0}^2 \subseteq \Kspace
  \quad\forall\, c_0 > 0
  \quad\Longrightarrow\quad
  \Kspace = L^2(\mathbb{R}),
\]
where the implication is by density of
$\bigcup_{c_0}\PW_{c_0}^2$ (Lemma~\ref{lem:pw-dense}).
The hypothesis $\PW_{c_0}^2 \subseteq \Kspace$
is a local completeness property in the inductive
limit of Paley--Wiener spaces.
This is structurally weaker in spectral support
than global completeness, but it is not a
reformulation: it is a local statement about
each finite bandwidth space that propagates to
global completeness only via the density step.
In particular, a hypothetical proof of H2
for a single value $c_0$ would not suffice;
the density argument requires all $c_0 > 0$.
\end{remark}

\begin{remark}[Role of Paper~VIII in Route~B]
\label{rem:paper8-in-B}
The nonvanishing estimates of Paper~VIII
(Lemma~\ref{lem:inherited}\ref{A:nonvanish}) enter
Route~B through H1: they prevent coefficient
collapse $a_n^{(c)} \to 0$ for
$n \le \kappa c$ by ensuring
$|\widehat{\Phi}_n^{(c)}| \ge c_1 c^{-1/2}$ on $E_c$,
which gives non-trivial inner products with
bandlimited functions. Making this precise is
part of Open Problem~\ref{prob:pw-transfer}.
\end{remark}

%%═══ 6. ROUTE C: CONDITIONAL COMPACTNESS
\section{Route~C: Conditional Compactness}
\label{sec:compact}

\begin{remark}[Compactness is a genuine additional hypothesis]
\label{rem:compact-not-free}
Since $\spc{\Hlim} \subseteq [0,1]$,
the operator $(\Hlim + I)^{-1}$ is well-defined and bounded
(Lemma~\ref{lem:inherited}\ref{A:bounded-sa}).
Compactness of $(\Hlim + I)^{-1}$ is strictly stronger:
it is not implied by boundedness, self-adjointness,
or SOT-convergence of $\Hc$ (Remark~\ref{rem:sot-limitation}).
\end{remark}

\begin{hypothesis}[Compact resolvent]
\label{hyp:compact-resolvent}
$(\Hlim + I)^{-1}$ is compact on $L^2(\mathbb{R})$.
\end{hypothesis}

\begin{proposition}[Conditional completeness via compact resolvent]
\label{prop:compact-route}
\textup{(Conditional on Hypothesis~\ref{hyp:compact-resolvent}.)}
If $(\Hlim + I)^{-1}$ is compact, then:
\begin{enumerate}[\rm(i)]
\item $\Hlim$ has purely discrete spectrum:
  $\spc{\Hlim} = \{\lambdainf{n}\}_{n \ge 0}$
  with finite multiplicities,
  accumulating only at $0$ or $1$;
\item $\{\Phiinf{n}\}$ is an orthonormal basis
  of $L^2(\mathbb{R})$;
\item $\Kspace = L^2(\mathbb{R})$.
\end{enumerate}
\end{proposition}
\begin{proof}
If $(\Hlim + I)^{-1}$ is compact, then by the spectral
theory of compact self-adjoint operators
\cite[Chapter~VI]{ReedSimon1975},
the spectrum of $(\Hlim + I)^{-1}$ consists of
eigenvalues accumulating only at $0$.
Translating back: $\Hlim = (\Hlim + I)^{-1} - I$
has purely discrete spectrum $\{\lambdainf{n}\}$
with finite multiplicities.
Since $\spc{\Hlim} \subseteq [0,1]$,
accumulation can only occur at $0$ or $1$.
The spectral theorem for bounded self-adjoint operators
with purely discrete spectrum then gives
the eigenfunction expansion in (ii)--(iii)
\cite[Theorem~VIII.6]{ReedSimon1975}.
\end{proof}

\begin{remark}[Mutual exclusivity with Paper~X, Open Problem~4]
\label{rem:op4-conflict}
Paper~X Open Problem~4 asks whether $\spc{\Hlim} = [0,1]$.
An operator with spectrum $[0,1]$ (purely continuous
or mixed spectrum filling an interval) cannot have
compact resolvent, since compact self-adjoint operators
have discrete spectrum with finite multiplicities.
Hence Hypothesis~\ref{hyp:compact-resolvent} and
$\spc{\Hlim} = [0,1]$ are mutually exclusive;
Paper~XIII does not resolve this tension.
\end{remark}

%%═══ 7. STRUCTURAL CLASSIFICATION THEOREM
\section{Structural Classification}
\label{sec:classification}

\begin{theorem}[Classification of completeness mechanisms]
\label{thm:classification}
The completeness problem $\Kspace = L^2(\mathbb{R})$
admits the following classification:
\begin{enumerate}[\rm(I)]
\item\label{class:A}
  \textup{(Mechanism~A: Equivalence reformulation.)}
  For the orthonormal system $\{\Phiinf{n}\}$,
  the Parseval identity and the frame lower bound
  are each \emph{equivalent} to completeness
  (Theorem~\ref{thm:route-A-equivalence}).
  They are structural characterisations, not
  independent proof routes.
  Frame-based arguments must be localised to
  individual $\PW_c^2$-spaces to avoid circularity.

\item\label{class:B}
  \textup{(Mechanism~B: Genuine analytic reduction.)}
  Under Hypotheses~\ref{hyp:coeff-stability}
  and~\ref{hyp:reconstruction},
  $\Kspace = L^2(\mathbb{R})$
  (Theorem~\ref{thm:main-B}).
  Route~B is not circular:
  it establishes the local containment
  $\PW_{c_0}^2 \subseteq \Kspace$ for all $c_0$
  and closes via density.
  The open core is H2
  (Hypothesis~\ref{hyp:reconstruction}).

\item\label{class:C}
  \textup{(Mechanism~C: Structural collapse.)}
  Under Hypothesis~\ref{hyp:compact-resolvent},
  $\{\Phiinf{n}\}$ is an orthonormal basis
  and $\spc{\Hlim}$ is purely discrete
  with finite multiplicities accumulating
  only at $0$ or $1$
  (Proposition~\ref{prop:compact-route}).
  This hypothesis is mutually exclusive with
  $\spc{\Hlim} = [0,1]$.
\end{enumerate}

\medskip
\begin{center}
\renewcommand{\arraystretch}{1.3}
\begin{tabular}{lllll}
\toprule
Route & Type & Circular? & Hypothesis & Open core \\
\midrule
A & Equivalence & Yes & None & ---
  (closed, negative) \\
B & Analytic reduction & No
  & H1+H2 & H2 (Hyp.~\ref{hyp:reconstruction}) \\
C & Structural collapse & No
  & Hyp.~\ref{hyp:compact-resolvent}
  & Spectral-type question \\
\bottomrule
\end{tabular}
\end{center}
\end{theorem}

\begin{remark}[Three contributions of Paper~XIII]
Paper~XIII makes three contributions to the series:
\begin{enumerate}[\rm(a)]
\item Theorem~\ref{thm:route-A-equivalence}:
  Route~A is an equivalence, not a reduction
  (a negative structural result that closes
  off a class of proof attempts);
\item Theorem~\ref{thm:main-B}:
  Route~B is the primary analytic route,
  with the open core precisely identified
  as H2;
\item The explicit identification of
  Hypotheses~\ref{hyp:coeff-stability},
  \ref{hyp:reconstruction},
  \ref{hyp:compact-resolvent}
  as the precise remaining obstructions.
\end{enumerate}
\end{remark}

%%═══ 8. OPEN PROBLEMS
\section{Open Problems}
\label{sec:open}

\begin{problem}[H1: coefficient stability in $\ell^2$]
\label{prob:coeff-stability}
Prove Hypothesis~\ref{hyp:coeff-stability}:
for $g \in \PW_{c_0}^2$,
$\ip{g}{\Phic{c}{n}} \to \ip{g}{\Phiinf{n}}$
in $\ell^2$ as $c \to \infty$.
This requires:
\begin{enumerate}[\rm(a)]
\item Strong eigenvector convergence for each
  fixed $n$ (Paper~XII, conditional on
  Gap-S + SOT-Faster);
\item Uniform tail control via
  Bessel: $\sum_{n>N}|\ip{g}{\Phic{c}{n}}|^2
  \le \norm{g}^2 - \sum_{n \le N}|\ip{g}{\Phic{c}{n}}|^2$.
\end{enumerate}
\end{problem}

\begin{problem}[H2: reconstruction in $L^2$]
\label{prob:pw-transfer}
Prove Hypothesis~\ref{hyp:reconstruction}:
for each $c_0 > 0$ and $g \in \PW_{c_0}^2$,
\[
  g = \sum_{n \ge 0} \ip{g}{\Phiinf{n}} \Phiinf{n}
  \quad \text{in } L^2(\mathbb{R}).
\]
This is the genuinely hard step:
it requires that $\Kspace \supseteq \PW_{c_0}^2$,
i.e.\ that $\{\Phiinf{n}\}$ spans $\PW_{c_0}^2$
(local completeness at bandwidth $c_0$).
It does not follow from H1 alone.
Note: by Theorem~\ref{thm:route-A-equivalence},
\emph{global} completeness $\Kspace = L^2(\mathbb{R})$
would imply H2, but H2 for all $c_0$ suffices
(via Lemma~\ref{lem:pw-dense}) to prove global completeness
without circular reasoning.
\end{problem}

\begin{problem}[Spectral type of $\Hlim$]
\label{prob:spectral-type}
Determine whether $\spc{\Hlim}$ is discrete,
continuous, or mixed:
\begin{enumerate}[\rm(a)]
\item Is $\spc{\Hlim} = \{\lambdainf{n}\}$?
\item Is $\spc{\Hlim} = [0,1]$
  (Paper~X, Open Problem~4)?
\item Is $(\Hlim + I)^{-1}$ compact
  (Hypothesis~\ref{hyp:compact-resolvent})?
\end{enumerate}
(b) and (c) are mutually exclusive.
\end{problem}

\begin{problem}[Operator identification]
\label{prob:bridge}
Prove $\Hlim = \overline{H_{\mathrm{spec}}}$,
where $H_{\mathrm{spec}}
:= \sum_n \lambdainf{n}\ip{\cdot}{\Phiinf{n}}\Phiinf{n}$.
Completeness of $\{\Phiinf{n}\}$ is necessary for
$H_{\mathrm{spec}}$ to be densely defined
(Paper~IX, Open Problem~7).
\end{problem}

%%═══ 9. STRUCTURAL SUMMARY
\section{Structural Summary}
\label{sec:summary}

Paper~XIII identifies and classifies three mechanisms
for the completeness problem $\Kspace = L^2(\mathbb{R})$:

\medskip
\begin{center}
\renewcommand{\arraystretch}{1.3}
\begin{tabular}{lllll}
\toprule
Route & Character & Open core & Blocks on \\
\midrule
A & Equivalence & ---
  & Circular (Thm.~\ref{thm:route-A-equivalence}) \\
B & Analytic reduction
  & H2 (Hyp.~\ref{hyp:reconstruction})
  & Local span of $\PW_c^2$ \\
C & Structural collapse
  & Spectral type
  & Hyp.~\ref{hyp:compact-resolvent}
  vs.~$\spc{\Hlim}=[0,1]$ \\
\bottomrule
\end{tabular}
\end{center}

\medskip
\noindent
Route~B is the only non-circular analytic route.
Its proof skeleton is complete;
the open core is H2 (local $\PW_c^2$-spanning),
grounded in Papers~VIII, IX, and~XII
without spectral-type assumptions on $\Hlim$.

%%%─── BIBLIOGRAPHY
\begin{thebibliography}{99}

\bibitem{PaperXIII-XI}
[Author],
\textit{Spectral structure for the limiting PSWF
operator}, Paper~XI, 2026.

\bibitem{PaperXIII-X}
[Author],
\textit{Mosco form convergence and the bridge
theorem}, Paper~X, 2026.

\bibitem{PaperXIII-IX}
[Author],
\textit{Conditional framework: domain, closability,
self-adjointness}, Paper~IX, 2026.

\bibitem{PaperXIII-VIII}
[Author],
\textit{Non-cancellation control and the sharp
$c^{-1/2}$ lower bound}, Paper~VIII, 2026.

\bibitem{PaperXIII-XII}
[Author],
\textit{A localization principle for spectral
projection stability}, Paper~XII, 2026.

\bibitem{Slepian1961}
D.~Slepian and H.~O.~Pollak,
\textit{Prolate spheroidal wave functions,
Fourier analysis and uncertainty~I},
Bell Syst.\ Tech.\ J.\ \textbf{40} (1961), 43--63.

\bibitem{Osipov2013}
A.~Osipov, V.~Rokhlin, H.~Xiao,
\textit{Prolate Spheroidal Wave Functions of
Order Zero}, Springer, 2013.

\bibitem{Kato1966}
T.~Kato,
\textit{Perturbation Theory for Linear Operators},
Springer, 1966.

\bibitem{ReedSimon1975}
M.~Reed and B.~Simon,
\textit{Methods of Modern Mathematical Physics,
Vol.~II}, Academic Press, 1975.

\bibitem{Conway1990}
J.~B.~Conway,
\textit{A Course in Functional Analysis},
2nd ed., Springer, 1990.

\bibitem{Rudin1991}
W.~Rudin,
\textit{Functional Analysis}, 2nd ed.,
McGraw-Hill, 1991.

\bibitem{Young2001}
R.~M.~Young,
\textit{An Introduction to Nonharmonic Fourier
Series}, Academic Press, 2001.

\end{thebibliography}

\end{document}
